\usepackage{xparse}    % 更灵活设置命令
\usepackage{mfirstuc}  % 大写首字母命令

% =========================================================
% Macro Library (AdvMacro II)
% 设计原则：
% 1) 保留原有命名接口，降低迁移成本
% 2) 按「大模块 -> 小模块」组织，便于查找与扩展
% 3) 基础宏在前，依赖宏在后，避免阅读时跳转
% =========================================================

\newcommand{\typolog}[2]{\noindent\ttt{#1}｜\xmuredbf{勘误}｜#2\par\vspace{0.1cm}}
\newcommand{\newlog}[2]{\noindent\ttt{#1}｜\fafugreenbf{完成}｜#2\par\vspace{0.1cm}}
\newcommand{\changelog}[2]{\noindent\ttt{#1}｜\xmubluebf{更新}｜#2\par\vspace{0.1cm}}


% 快速记号
\newcommand{\ykbar}{\dfrac{\tilde\delta}{\alpha\,\beta}}
\newcommand{\ckbar}{\bka{\dfrac{\tilde\delta}{\alpha\,\beta}-\delta}}

% =========================================================
% A. 基础快捷层 (Base Shortcuts)
% =========================================================

% A1. 文本包装与上下角标快捷
\newcommand{\n}[1]{\text{#1}}
\newcommand{\w}[1]{^{#1}}
\newcommand{\s}[1]{_{#1}}
\newcommand{\wn}[1]{^{\text{#1}}}
\newcommand{\sn}[1]{_{\text{#1}}}

\newcommand{\wsii}{^{\infty}_{i=1}}
\newcommand{\wsni}{^{n}_{i=1}}
\newcommand{\wsito}{^{\infty}_{t=0}}

% A2. 常用显示模式/短语
\newcommand{\dis}{\displaystyle}
\newcommand{\st}{\text{s.t.}}
\newcommand{\iid}{\text{i.i.d.}}


% =========================================================
% B. 字体与字母体系 (Fonts & Alphabets)
% =========================================================

% B1. 黑板体、花体、草体
\newcommand{\bbrn}{\bbr[n]}

\NewDocumentCommand{\bbr}{o}{%
  \mathbb{R}\IfValueT{#1}{^{#1}}%
}

\NewDocumentCommand{\bb}{m o}{%
  \mathbb{#1}\IfValueT{#2}{^{#2}}%
}

\RenewDocumentCommand{\cal}{m o}{%
  \mathcal{#1}\IfValueT{#2}{^{#2}}%
}

\NewDocumentCommand{\scr}{m o}{%
  \mathscr{#1}\IfValueT{#2}{^{#2}}%
}

% B2. 常用目标函数符号
\newcommand*{\lagrangian}{\mathcal{L}}


% =========================================================
% C. 颜色体系 (Color Helpers)
% =========================================================

\newcommand{\xmured}[1]{{\color[HTML]{D7422A}#1}}
\newcommand{\xmublue}[1]{{\color[HTML]{003C88}#1}}
\newcommand{\fafugreen}[1]{{\color[HTML]{00833A}#1}}

\newcommand{\noxmubluebf}[1]{\noindent\textcolor[HTML]{003C88}{\textbf{#1}}}
\newcommand{\xmuredbf}[1]{\textcolor[HTML]{D7422A}{\textbf{#1}}}
\newcommand{\xmubluebf}[1]{\textcolor[HTML]{003C88}{\textbf{#1}}}
\newcommand{\fafugreenbf}[1]{\textcolor[HTML]{00833A}{\textbf{#1}}}


% =========================================================
% D. 间距体系 (Spacing System)
% =========================================================

% D1. 文档垂直间距：v = 增加, m = 减少
\newcommand{\spacexs}{\vspace{1pt}}
\newcommand{\spaces}{\vspace{2pt}}
\newcommand{\spacem}{\vspace{4pt}}
\newcommand{\spacel}{\vspace{6pt}}
\newcommand{\mspacexs}{\vspace{-1pt}}
\newcommand{\mspaces}{\vspace{-2pt}}
\newcommand{\mspacem}{\vspace{-4pt}}
\newcommand{\mspacel}{\vspace{-6pt}}

% D2. 数学水平间距
\newcommand{\mathspacexs}{\mspace{1mu}}
\newcommand{\mathspaces}{\mspace{2mu}}
\newcommand{\mathspacem}{\mspace{4mu}}
\newcommand{\mathspacel}{\mspace{6mu}}

% D3. 统一调整数学模式中的 \, 宽度（正文模式保持原行为）
\let\originalthinspace\,
\renewcommand{\,}{\ifmmode \mathspaces \else \originalthinspace \fi}



% =========================================================
% E. 括号与分隔符体系 (Delimiters)
% =========================================================

% E1. 基础自适应括号（bk 系列保留 left/right）
\newcommand{\bka}[1]{\left(\mathspacexs#1\mathspacexs\right)}
\newcommand{\bks}[1]{\left[\mathspacexs#1\mathspacexs\right]}
\newcommand{\bkd}[1]{\left\{\mathspacexs#1\mathspacexs\right\}}
\newcommand{\bkf}[1]{\left|\mathspacexs#1\mathspacexs\right|}
\newcommand{\bkg}[1]{\left\|\mathspacexs#1\mathspacexs\right\|}
\newcommand{\bkas}[1]{\left(\mathspacexs#1\mathspacexs\right]}
\newcommand{\bksa}[1]{\left[\mathspacexs#1\mathspacexs\right)}

% E2. 手动尺寸括号快捷层
% 约定：
% q  = big
% p  = Big
% qp = bigg
% pp = Bigg
% 括号外侧额外保留一点点小间隙
\newcommand{\delimout}{\mathspaces}

\newcommand{\pa}[1]{\delimout\bigl(\mathspacexs#1\mathspacexs\bigr)\delimout}
\newcommand{\ps}[1]{\delimout\bigl[\mathspacexs#1\mathspacexs\bigr]\delimout}
\newcommand{\pd}[1]{\delimout\bigl\{\mathspacexs#1\mathspacexs\bigr\}\delimout}
\newcommand{\pf}[1]{\delimout\bigl|\mathspacexs#1\mathspacexs\bigr|\delimout}
\newcommand{\pg}[1]{\delimout\bigl\|\mathspacexs#1\mathspacexs\bigr\|\delimout}

\newcommand{\qpa}[1]{\delimout\Bigl(\mathspacexs#1\mathspacexs\Bigr)\delimout}
\newcommand{\qps}[1]{\delimout\Bigl[\mathspacexs#1\mathspacexs\Bigr]\delimout}
\newcommand{\qpd}[1]{\delimout\Bigl\{\mathspacexs#1\mathspacexs\Bigr\}\delimout}
\newcommand{\qpf}[1]{\delimout\Bigl|\mathspacexs#1\mathspacexs\Bigr|\delimout}
\newcommand{\qpg}[1]{\delimout\Bigl\|\mathspacexs#1\mathspacexs\Bigr\|\delimout}

\newcommand{\ppa}[1]{\delimout\biggl(\mathspacexs#1\mathspacexs\biggr)\delimout}
\newcommand{\pps}[1]{\delimout\biggl[\mathspacexs#1\mathspacexs\biggr]\delimout}
\newcommand{\ppd}[1]{\delimout\biggl\{\mathspacexs#1\mathspacexs\biggr\}\delimout}
\newcommand{\ppf}[1]{\delimout\biggl|\mathspacexs#1\mathspacexs\biggr|\delimout}
\newcommand{\ppg}[1]{\delimout\biggl\|\mathspacexs#1\mathspacexs\biggr\|\delimout}

\newcommand{\qa}[1]{\delimout(\mathspacexs#1\mathspacexs)\delimout}
\newcommand{\qs}[1]{\delimout[\mathspacexs#1\mathspacexs]\delimout}
\newcommand{\qd}[1]{\delimout\{\mathspacexs#1\mathspacexs\}\delimout}
\newcommand{\qf}[1]{\delimout|\mathspacexs#1\mathspacexs|\delimout}
\newcommand{\qg}[1]{\delimout\|\mathspacexs#1\mathspacexs\|\delimout}

% E3. 竖线和斜杠分隔符
\newcommand{\islash}{{\mathspaces}/{\mathspaces}}
\newcommand{\bigslash}{{\mathspaces}\big/{\mathspaces}}
\newcommand{\biggslash}{{\mathspaces}\bigg/{\mathspaces}}
\newcommand{\Bigslash}{{\mathspaces}\Big/{\mathspaces}}
\newcommand{\Biggslash}{{\mathspaces}\Bigg/{\mathspaces}}

\newcommand{\ivbar}{{\mathspaces}|{\mathspaces}}
\newcommand{\bigbar}[1]{{\mathspaces}\big|_{\,#1}}
\newcommand{\biggbar}[1]{{\mathspaces}\bigg|_{\,#1}}
\newcommand{\Bigbar}[1]{{\mathspaces}\Big|_{\,#1}}
\newcommand{\Biggbar}[1]{{\mathspaces}\Bigg|_{\,#1}}

% E4. 矩阵/向量括号模板：默认给一个固定尺寸，不再自动伸缩
\newcommand{\matspace}{\mspace{2mu}}

\newcommand{\matbka}[1]{\left(\matspace\begin{matrix}#1\end{matrix}\matspace\right)}
\newcommand{\matbks}[1]{\left[\matspace\begin{matrix}#1\end{matrix}\matspace\right]}
\newcommand{\matbkf}[1]{\left|\matspace\begin{matrix}#1\end{matrix}\matspace\right|}
\newcommand{\vecbka}[1]{\left(\matspace\begin{matrix}#1\end{matrix}\matspace\right)}
\newcommand{\vecbks}[1]{\left[\matspace\begin{matrix}#1\end{matrix}\matspace\right]}

% E5. 常用加修饰括号（bar/hat/tilde + 下标文本）
\newcommand{\qab}[2]{\bar{\bka{#1}}}
\newcommand{\qah}[2]{\hat{\bka{#1}}}
\newcommand{\qat}[2]{\tilde{\bka{#1}}}
\newcommand{\qabw}[2]{\qa{\mathspacexs\bar{#1}^{\text{#2}}\mathspacexs}}
\newcommand{\qahw}[2]{\qa{\mathspacexs\hat{#1}^{\text{#2}}\mathspacexs}}
\newcommand{\qatw}[2]{\qa{\mathspacexs\tilde{#1}^{\text{#2}}\mathspacexs}}
\newcommand{\qabs}[2]{\qa{\mathspacexs\bar{#1}_{\text{#2}}\mathspacexs}}
\newcommand{\qahs}[2]{\qa{\mathspacexs\hat{#1}_{\text{#2}}\mathspacexs}}
\newcommand{\qats}[2]{\qa{\mathspacexs\tilde{#1}_{\text{#2}}\mathspacexs}}
\newcommand{\qabws}[3]{\qa{\mathspacexs\bar{#1}^{\text{#2}}_{\text{#3}}\mathspacexs}}
\newcommand{\qahws}[3]{\qa{\mathspacexs\hat{#1}^{\text{#2}}_{\text{#3}}\mathspacexs}}
\newcommand{\qatws}[3]{\qa{\mathspacexs\tilde{#1}^{\text{#2}}_{\text{#3}}\mathspacexs}}


% =========================================================
% F. 运算与统计算子 (Operations & Statistical Operators)
% =========================================================

% F1. 极限与收敛相关
\newcommand{\nti}{n\to\infty}
\newcommand{\limn}{\lim_{n\to\infty}}

% F2. 行内分式与根号格式
\newcommand{\ifrac}[2]{%
  \mathchoice
    {\frac{#1}{#2}}
    {\scaleobj{1.3}{\frac{#1}{#2}}}
    {\frac{#1}{#2}}
    {\frac{#1}{#2}}
}

\makeatletter
\NewDocumentCommand{\dsq}{o}{\@ifnextchar\bgroup{\dsq@a[#1]}{\@ifnextchar- {\dsq@neg[#1]}{\dsq@one[#1]}}}
\def\dsq@a[#1]#2{\IfValueTF{#1}{\dis\sqrt[#1]{\,#2}}{\dis\sqrt{\,#2}}}
\def\dsq@neg[#1]-#2{\dsq@a[#1]{-#2}}
\def\dsq@one[#1]#2{\dsq@a[#1]{#2}}
\makeatother

% F3. 统计算子：只保留裸算子，不再绑定括号类型
\newcommand{\mean}{\text{E}}
\newcommand{\meant}{\text{E}_t}
\newcommand{\meano}{\text{E}_0}
\newcommand{\var}{\text{Var}}
\newcommand{\cov}{\text{Cov}}
\newcommand{\avar}{\text{avar}}
\newcommand{\se}{\text{se}}
\newcommand{\pr}{\Pr}

% F4. 极值与集合上下界显示
\renewcommand{\max}{\mathop{\text{max}}}
\renewcommand{\min}{\mathop{\text{min}}}

\newcommand{\maxlim}{\mathop{\text{max}}}
\newcommand{\minlim}{\mathop{\text{min}}}
\newcommand{\suplim}{\mathop{\text{sup}}}
\newcommand{\inflim}{\mathop{\text{inf}}}
\newcommand{\cuplim}{\bigcup\limits}
\newcommand{\caplim}{\bigcap\limits}


% =========================================================
% G. 微积分与差分体系 (Calculus & Difference)
% =========================================================

% G1. 微分/偏导/差分符号
\renewcommand{\d}{\mathrm{d}\mathspaces}
\newcommand{\dd}{\mathrm{d}^2\mathspaces}
\newcommand{\p}{\partial\mathspaces}
\newcommand{\pp}{\partial^2\mathspaces}
\newcommand{\D}{\Delta\mathspaces}
\newcommand{\DD}{\Delta^2\mathspaces}

% G2. 导数模板
\newcommand{\dfracd}[1]{\frac{\mathrm{d}}{\d{#1}}}
\newcommand{\dfracp}[1]{\frac{\partial}{\p{#1}}}
\newcommand{\dfracdd}[1]{\frac{\mathrm{d}^2}{\d{#1}^2}}
\newcommand{\dfracpp}[1]{\frac{\partial^2}{\p{#1}^2}}


% =========================================================
% H. 箭头与上下角标体系 (Arrows & Scripts)
% =========================================================

% H1. 概率论常用收敛箭头
\newcommand{\pto}{\xrightarrow{p}}
\newcommand{\dto}{\xrightarrow{d}}
\newcommand{\asto}{\xrightarrow{a.s.}}
\newcommand{\wto}{\xrightarrow{w}}

% H2. 通用上下角标后缀
\newcommand{\inv}{^{-1}}
\newcommand{\tra}{^{\top}}
\newcommand{\nt}{_{t+1}}
\newcommand{\nnt}{_{t+2}}
\newcommand{\lt}{_{t-1}}
\newcommand{\llt}{_{t-2}}


% =========================================================
% I. 符号风格与分布名 (Styling & Distributions)
% =========================================================

% I1. 其他通用符号
\newcommand{\abfun}{\,\cdot\,}
\renewcommand{\i}{\infty}

% I2. 加粗数学符号（向量/矩阵）
\newcommand{\bs}[1]{\symbfit{#1}}
\newcommand{\bbs}[1]{\bar{\symbfit{#1}}}
\newcommand{\hbs}[1]{\hat{\symbfit{#1}}}
\newcommand{\tbs}[1]{\tilde{\symbfit{#1}}}

\newcommand{\bsp}{\symbfit{p}}
\newcommand{\bsx}{\symbfit{x}}
\newcommand{\bsy}{\symbfit{y}}
\newcommand{\bsz}{\symbfit{z}}
\newcommand{\bsA}{\symbfit{A}}
\newcommand{\bsM}{\symbfit{M}}
\newcommand{\bsI}{\symbfit{I}}
\newcommand{\bsX}{\symbfit{X}}

% I3. 常见概率分布名
\newcommand{\poisson}{\text{Poisson}}
\newcommand{\normal}{\text{Normal}}
\newcommand{\exponential}{\text{Exponential}}
\newcommand{\bernoulli}{\text{Bernoulli}}
\newcommand{\binomial}{\text{Binomial}}
\newcommand{\geometric}{\text{Geometric}}
\newcommand{\uniform}{\text{Uniform}}
\newcommand{\lognormal}{\text{Lognormal}}
\newcommand{\cauchy}{\text{Cauchy}}
\newcommand{\negabinomial}{\text{NBinomial}}
\newcommand{\gam}{\text{Gamma}}


% =========================================================
% J. 高层快捷模板 (High-level Shorthands)
% =========================================================

% J1. 单字母括号快捷
\newcommand{\bkk}{\bka{k}}
\newcommand{\bkt}{\bka{t}}
\newcommand{\bku}{\bka{u}}
\newcommand{\bkv}{\bka{v}}
\newcommand{\bkx}{\bka{x}}
\newcommand{\bki}{\bka{i}}
\newcommand{\bkj}{\bka{j}}
\newcommand{\bky}{\bka{y}}
\newcommand{\bkz}{\bka{z}}
\newcommand{\bkA}{\bka{A}}
\newcommand{\bkX}{\bka{X}}
\newcommand{\bko}{\bka{0}}

% J2. 单字母加粗括号快捷
\newcommand{\bkbsp}{\bka{\symbfit{p}}}
\newcommand{\bkbsx}{\bka{\symbfit{x}}}
\newcommand{\bkbsy}{\bka{\symbfit{y}}}
\newcommand{\bkbs}[1]{\bka{\symbfit{#1}}}

% J3. 常见时序/样本表达快捷
\newcommand{\bkct}{\bka{c_t}}
\newcommand{\bkcnt}{\bka{c_{t+1}}}
\newcommand{\bkkt}{\bka{k_t}}
\newcommand{\bkknt}{\bka{k_{t+1}}}
\newcommand{\bkkstar}{\bka{k^*}}
\newcommand{\bkdxn}{\bkd{x_n}}
\newcommand{\bkdxt}{\bkd{x_t}}
\newcommand{\bkdyt}{\bkd{y_t}}
\newcommand{\bkdut}{\bkd{u_t}}
\newcommand{\bkXYi}{\bkd{X_i,Y_i}}

% J4. 多变量组合快捷
\newcommand{\bkab}{\bka{a,b}}
\newcommand{\bkxy}{\bka{x,y}}
\newcommand{\bkxz}{\bka{x,z}}
\newcommand{\bkyz}{\bka{y,z}}
\newcommand{\bkxyz}{\bka{x,y,z}}
\newcommand{\bkXY}{\bka{X,Y}}
\newcommand{\bkKL}{\bka{K,L}}
\newcommand{\bkKAL}{\bka{K,AL}}
\newcommand{\bkoo}{\bka{0,0}}
\newcommand{\bkdXY}{\bkd{X,Y}}
\newcommand{\bkit}{\bka{i,t}}
\newcommand{\bkqq}{\bka{q_1,q_2}}
\newcommand{\bkxx}{\bka{x_1,x_2}}
\newcommand{\bkXX}{\bka{X_1,X_2}}
\newcommand{\bkyy}{\bka{y_1,y_2}}

% J5. 序列书写模板
\newcommand{\xx}[1]{#1_1,#1_2}
\newcommand{\xxstar}[1]{#1^*_1,#1^*_2}
\newcommand{\xxprime}[1]{#1'_1,#1'_2}
\newcommand{\xxx}[1]{#1_1,#1_2,#1_3}
\newcommand{\sexn}[1]{#1_1,\dots,#1_n}
\newcommand{\sexs}[1]{#1_1,\dots,#1_s}
\newcommand{\sexxn}[1]{#1_1,#1_2,\dots,#1_n}
\newcommand{\sexxd}[1]{#1_1,#1_2,\dots}

% J6. 增长率模板
\newcommand{\grt}[1]{\dfrac{\dot #1}{#1}}
\newcommand{\grta}[1]{\dfrac{\dot{\bka{#1}}}{#1}}
\newcommand{\grts}[1]{\dfrac{\dot{\bks{#1}}}{#1}}
\newcommand{\grtt}[1]{\dfrac{\dot #1\bkt}{#1\bkt}}
\newcommand{\ingrt}[1]{\dot{#1}\islash{#1}}
\newcommand{\ingrta}[1]{\dot{\bka{#1}}\islash{#1}}
\newcommand{\ingrts}[1]{\dot{\bks{#1}}\islash{#1}}

\newcommand{\dfracb}[2]{\dfrac{\bar{#1}}{\bar{#2}}}


% ==========================================
% ==========================================
% config/mathcmds.tex — 数学符号与算子宏
% ==========================================
% ==========================================

% ---------- 概率与统计 ----------
\DeclareMathOperator{\Var}{Var}        % 方差
\DeclareMathOperator{\Cov}{Cov}        % 协方差
\DeclareMathOperator{\Corr}{Corr}      % 相关系数
\DeclareMathOperator{\SD}{SD}          % 标准差
\DeclareMathOperator{\SE}{se}          % 标准误

% ---------- 优化 ----------
\DeclareMathOperator*{\argmax}{arg\,max}
\DeclareMathOperator*{\argmin}{arg\,min}
\DeclareMathOperator*{\maximize}{maximize}
\DeclareMathOperator*{\minimize}{minimize}

% ---------- 括号与范数 ----------
\DeclarePairedDelimiter{\abs}{\lvert}{\rvert}
\DeclarePairedDelimiter{\norm}{\lVert}{\rVert}
\DeclarePairedDelimiter{\inp}{\langle}{\rangle}      % 内积
\DeclarePairedDelimiter{\floor}{\lfloor}{\rfloor}
\DeclarePairedDelimiter{\ceil}{\lceil}{\rceil}

% ---------- 对数 ----------
\let\oldln\ln
\RenewDocumentCommand{\ln}{}{\operatorname{ln}}
\let\oldlog\log
\RenewDocumentCommand{\log}{}{\operatorname{log}}
\DeclareMathOperator{\tr}{tr}


\newcommand{\means}{\quad\Rightarrow\quad}
\newcommand{\rmeans}{\Rightarrow\quad}